\section{The Hunger of Charity}

The passage from justice to charity becomes especially delicate at the Eucharist. Christ's command is not a metaphor awaiting a spiritual substitute. ``Unless you eat the flesh of the Son of man and drink his blood, you have no life in you.'' He then gives the promise: ``He who eats my flesh and drinks my blood abides in me, and I in him,'' and, ``He who eats me will live because of me.''\endnote{John 6:53--57; CCC 1382--1392. John 6 is read within the Church's sacramental doctrine.}

There is no shame in beginning with duty. The Church's precepts give objective shape to a creature whose desire is inconsistent. One who attends Mass or receives according to the Church's law before feeling hunger is no hypocrite. Obedience can precede delight, just as a person may keep faith with someone before recovering the affection that once made fidelity easy. But duty is meant to guard an encounter. If it becomes evidence that the account is settled, the sacrament itself is being used against its end.

Nor does desire mature by leaving rite behind for an allegedly interior union. The sacramental sign is not a shell from which the mystic escapes. Worthy reception gives Christ himself; interior hunger moves toward the gift he chose. When grave sin has broken communion, charity does not rename presumption as intimacy: it seeks Reconciliation before receiving.\endnote{CCC 1382--1389, 1395. Questions of disposition, especially under scrupulosity, belong with a regular confessor.}

Aquinas calls the Eucharist the sacrament of charity because it perfects union with Christ who suffered. Its fruits are not passing emotions: as spiritual food it preserves and increases grace and rekindles charity; its principal fruit is intimate union with Christ, inseparable from his members.\endnote{ST III, q.~73, a.~3; q.~79, aa.~1 and 4--8; CCC 1391--1397.}

Benedict XVI draws out the strangeness of this food. Ordinary food is assimilated into the eater; here Christ ``assimilates us'' to himself. This conformity does not erase individuality; it opens the person beyond egocentrism into Christ's Body and therefore toward the difficult neighbor.\endnote{Benedict XVI, Corpus Christi homily, 23 June 2011, drawing on Augustine, \emph{Confessions} VII.10.16.}

Thérèse of Lisieux expresses the same priority. She places before her longing Christ's desire to unite himself to us and dwell within. Human desire is not passive but freed from self-display. The soul hungers because it has first been sought; even its capacity to receive is gift.\endnote{\emph{C'est la confiance}, nos.~18--22, especially 22, receiving Thérèse's \emph{Act of Oblation}.}

Here duty ripens into desire without being discarded. The communicant no longer asks only, \emph{What must I do?} but, \emph{How could I neglect the One who gives himself as food?} Yet ``The Eucharist commits us to the poor.''\endnote{CCC 1397: Eucharistic union is tested in mercy and justice toward the visible neighbor.} To receive the Body while refusing his suffering members contradicts the ``Amen'' spoken at the altar.

The hunger of charity is therefore not hunger for a religious experience. It is desire for Christ---his truth, his commandments, his sacramental gift, his members, his will---because he is Christ.
